\chapter{Design decisions and roads not taken}
\label{ch:decisions}

A library is shaped as much by what it leaves out as by what it includes. Each item below was
tried or seriously considered, and each has a concrete reason for its current status. The full
records are in \file{docs/MULTIPLICATION.md}, \file{docs/DIVISION.md} and
\file{docs/STRING\_CONVERSION.md}.

\section{Principles that recur}

\begin{enumerate}
  \item \textbf{Measure end to end, on the shapes that really occur.} Microbenchmarks of one
        algorithm in isolation repeatedly gave wrong crossovers: they missed setup costs, the
        pool's sleep state, or lopsided shapes. The thresholds come from timing the dispatcher's
        actual choices over grids of sizes and shapes.
  \item \textbf{Re-measure when a building block changes.} Faster kernels moved the Karatsuba
        leaf from 48 to 24 limbs; faster NTT butterflies retired the Toom-3 band and moved the
        NTT entry from 2560 to 320 limbs per operand. A threshold is only correct for the code it
        was measured with.
  \item \textbf{Change the dependency structure, not the instruction count.} The biggest
        small-size gains (two carry chains, fused NTT layers, interleaved twiddle tables, fused
        pointwise products) came from shortening dependency chains and memory passes.
  \item \textbf{Build big algorithms on the dispatcher.} Division and conversion ask
        \code{Multiply} for their products, so every multiplication improvement propagates.
  \item \textbf{Silent slow paths are bugs.} A correct-but-quadratic fallback is invisible to
        correctness tests, so fallbacks are counted and tested for (\cref{sec:newton}).
\end{enumerate}

\section{Rejected or retired}

\begin{description}[leftmargin=0pt,style=nextline,itemsep=6pt]
  \item[Schönhage--Strassen multiplication]
        A parameter sweep found no winning size band for the single-level implementation
        considered here. Each three-prime NTT block accepts up to $2^{26}$ coefficients,
        about 323 million digits per balanced operand. Further convolution residues of that
        length extend one product to about 1.79 billion digits per operand before block
        accumulation. Factoring the cross-twiddles and
        distributing matrix tiles across the pool reduced its cost within that range
        (\cref{sec:mfa,sec:nttthreads}). Measured times appear in \cref{ch:practice}.
  \item[Toom-3 as a top-level band]
        Once a separate band between Karatsuba and NTT (total 2560--5120 limbs). The vectorized NTT
        overtook it everywhere, so the band was retired. Toom-3 now lives inside the Karatsuba
        recursion instead (\cref{sec:toom}).
  \item[Toom-4 and Toom-5]
        Toom-4 is implemented but off: no gain on the M1 before the NTT starts. Toom-5 was
        3 times slower than Karatsuba at its first active size (512 limbs): its evaluation and
        interpolation cost more than the multiplications they save.
  \item[Two-row unrolled schoolbook]
        Measured flat (\cref{sec:flagchain}): it kept the limb-sized carry chain.
  \item[Barrett reduction for the NTT primes]
        Slower than the compiler's own code for division by a constant in this loop. Shoup's
        method (precomputed per twiddle) replaced both.
  \item[Montgomery reduction for the Goldilocks prime]
        That prime already has a special-form reduction of a few additions; Montgomery would only
        add conversions.
  \item[Short (high-half-only) products in Newton division]
        Implemented exactly and measured flat. At the sizes where Newton runs, products go through
        the NTT, where computing ``half a product'' as two smaller products costs about the same as
        one full product. Shortening the \emph{transform} (wrap-around products) worked instead
        (\cref{sec:wrap}).
  \item[Reusing a prepared transform inside Newton chunks]
        Cut the number of transforms but not the wall-clock time: the three primes already run in
        parallel, so fewer transforms did not shorten the critical path.
  \item[Splitting decimal conversion at powers $10^{18 \cdot 2^i}$]
        Unbalanced splits made it slower than the quadratic method (\cref{ch:conv}).
  \item[Per-layer threading of the NTT]
        Starting threads for every butterfly layer spent most of its time in synchronization.
        Coarse units (whole transforms, or large row ranges) replaced it.
  \item[Profile-guided optimization]
        No measurable change: the hot loops are already fully inlined and unrolled.
  \item[Storing numbers in base $10^{18}$]
        Would make printing trivial and all arithmetic slower; arithmetic dominates real work.
\end{description}

\section{Open opportunities}

The project's notes name these as the most promising remaining directions:

\begin{itemize}
  \item a half-size reciprocal with two quotient blocks for divisions of ratio about 2 (the shape
        of GMP's \code{mu\_div});
  \item block-recursive (Mulders) division and Schönhage's subquadratic GCD, if a workload needs
        them;
  \item removing the power-of-two staircase above $2^{20}$ coefficients in threaded builds.
\end{itemize}
